\subsection{Effiziente Terminal-Node-Evaluierung}

Ein weiterer Beschleunigungsschritt betrifft die Evaluierung von Terminal Nodes.
Eine naive Herangehensweise würde für jedes Paar von Informationssets die Utility berechnen, was eine $O(n^2)$ Komplexität bedeutet, wobei $n$ die Anzahl der Informationssets pro Spieler ist.

In Pokerspielen kann diese Komplexität auf $O(n)$ reduziert werden, indem die Informationssets nach ihrer Handstärke sortiert werden.
Da die Utility nur von der relativen Rangordnung der Hände abhängt, kann die Evaluierung durch geschickte Indizierung effizient durchgeführt werden.
Für abhängige Verteilungen, wie sie in Poker auftreten, wird das Inklusions-Exklusions-Prinzip verwendet, um die korrekte Berechnung zu gewährleisten.
